% !TEX root=../manuscript.tex
\section{Evaluation \label{sec:evaluation}}
We design our evaluation to answer the following research questions.
\begin{itemize}
    \item \textbf{RQ1:} How does \mainname compare with its base models across model scales and benchmarks?
    \item \textbf{RQ2:} How does each component of \mainname contribute to the overall performance, and what runtime cost does it introduce?
    \item \textbf{RQ3:} How does \mainname compare with other reliability-oriented generation approaches, including rejection sampling, token-level constrained decoding, completion-engine repair, and iterative compiler-feedback repair?
    \item \textbf{RQ4:} How does \mainname compare with variants that generate the type derivation tree and the program separately instead of simultaneously?
\end{itemize}

\subsection{Experimental Setup}

\subsubsection{Benchmark}
To evaluate the generality of \mainname, our evaluation considers two languages with very different characteristics, SuFu~\cite{DBLP:journals/pacmpl/JiZPXH24} and Java.
%Correspondingly, we used two datasets: the SuFu dataset and the MBJP~\cite{DBLP:conf/iclr/AthiwaratkunGWL23} dataset.
\begin{itemize}
 \item \textbf{SuFu} is an ML-style functional language designed for automatic program optimization. This language uses a customized type system for extracting operations for optimization on data structures.
 %
 This language has limited code resources, making it representative for evaluating \mainname under low-resource conditions.

 We collect SuFu programs from its GitHub repository, which comprises 290 programs with test cases. For each program, we manually supply a natural language description, along with the relevant inductive definitions and library functions, thus collecting a dataset of code generation tasks.
 %
 This dataset is then split randomly into a training set of 80\% tasks and a test set of 20\% tasks.
 
 \item \textbf{Java} is a popular imperative language with a sophisticated type system. In a balance between the implementation cost and the expressiveness, we implement a subset of Java that preserves its core syntax while omitting advanced features like lambda expressions.

 We consider the MBJP dataset~\cite{DBLP:conf/iclr/AthiwaratkunGWL23} in our evaluation, which is the Java version of the popular MBPP dataset. Our Java subset covers the MBJP tasks admitted by our toolchain, which we split randomly into a training set of 608 tasks and a held-out test set of 67 tasks.

 \item \textbf{HumanEval-Java} and \textbf{TransCoder-GFG} are two additional Java benchmarks that broaden the evaluation beyond MBJP. HumanEval-Java adapts the HumanEval problems to Java~\cite{DBLP:journals/corr/abs-2107-03374}, and TransCoder-GFG collects the Java programs released with TransCoder~\cite{DBLP:conf/nips/LachauxROC20}; both follow the task format of MBJP. 
HumanEval-Java uses a fixed split of 146 training and 16 test tasks, and TransCoder-GFG uses a fixed split of 414 training and 103 test tasks.
\end{itemize}

More information about SuFu, Java, and the dataset construction can be found in Appendix~\ref{appendix-benchmark}.

\subsubsection{Models}
We implement \mainname based on CodeT5~\cite{DBLP:conf/emnlp/0034WJH21}, a family of encoder-decoder models for code understanding and generation. We adopt the 220M variant of CodeT5 as our base model for subsequent training and refer to our model as \mainname-220M. 
We also evaluate \mainname on T5Gemma2~\cite{zhang2025t5gemma2seeingreading}, a more recent encoder-decoder model family built upon Gemma 3. We adopt the 2B variant of T5Gemma2 as our base model and refer to this model as \mainname-2B. 

\subsubsection{Evaluation Metrics} 
\label{sec:evaluation-metrics}
To evaluate the quality of code generation, we adopt the metrics below, which respectively measure the correctness, effectiveness, and type correctness:

\begin{itemize}
    \item \textbf{Pass@k (k=1,10)~\cite{DBLP:journals/corr/abs-2107-03374}}: These metrics quantify the proportion of problems where at least one of the top-k generated candidates passes all reference test cases. Each generated candidate is executed against the test cases provided in the benchmark, and a candidate is considered correct only if it produces the expected output for all test cases. 
    %It's the standard and widely used method to quantify the correctness of generated code.

    \item \textbf{First Success Position (FSP)}: This metric is the mean rank of the first functionally correct candidate among the generated ones across the test set. It measures how early in the candidate list a correct solution appears, assessing the efficiency of output ranking; a problem with no correct candidate contributes the candidate-list length (10).

    \item \textbf{Compilation Error Rate (CER)}: This metric quantifies the proportion of generated candidates that fail to compile. 
    Since \mainname generates programs with type correctness guarantees, CER reflects the type system's effectiveness in catching compile-time errors. 
\end{itemize}

\subsubsection{Implementation Details}

We fine-tuned both \mainname-220M (based on CodeT5-220M) and \mainname-2B (based on T5Gemma2-2B) following a standard supervised fine-tuning procedure with the same set of hyperparameters to ensure a fair comparison. Each model was fine-tuned on the training set until convergence.
Since \mainname introduces a new representation that differs from standard code tokens and thus faces a cold-start problem, we first perform an initial training phase using a supervised fine-tuning dataset from OpenCoder~\cite{DBLP:journals/corr/abs-2411-04905} to adapt the model to this representation.
At the same scale and on the same benchmark, each baseline--\mainname{} pair is fine-tuned on the same training set. The exception is the two additional Java benchmarks, HumanEval-Java and TransCoder-GFG. There, \mainname-2B builds on its MBJP fine-tuning, which avoids the cold-start problem, and is further trained on the training splits of the two benchmarks. The baselines, which generate ordinary code tokens and face no cold-start problem, are instead trained from the original T5Gemma2-2B, with the MBJP training tasks added to their own training sets for joint training.

\subsection{Experimental Results}

\subsubsection{RQ1: Performance of \mainname}
\autoref{tab:model-results} summarizes the overall performance of \mainname, with models grouped by scale.
For each model scale and benchmark, we compare the fine-tuned baseline (CodeT5-220M or T5Gemma2-2B) with our \mainname{} model on the metrics above; the 2B group additionally covers the two Java benchmarks, HumanEval-Java and TransCoder-GFG.

\begin{table}[ht]
\centering
\begin{threeparttable}
\caption{Comparison of Accuracy and Compilation Error Rate between Base Models and \mainname at the 220M and 2B Scales}
\label{tab:model-results}
\begin{tabular}{
    l 
    l
    l
    S[table-format=2.2] 
    S[table-format=2.2] 
    S[table-format=1.2] 
    S[table-format=2.2]
}
\toprule
Scale & Benchmark & Model & {pass@1 (\%)\tnote{\updir}} & {pass@10 (\%)\tnote{\updir}} & {FSP\tnote{\downdir}} & {CER (\%)\tnote{\downdir}} \\
\midrule
\multirow{4}{*}{220M}
        & SuFu & CodeT5-220M         & 24.14 & 32.76 & 7.03 & 83.10 \\
        &      & \mainname-220M      & \textbf{37.93} & \textbf{53.45} & \textbf{5.03} & \textbf{0.00} \\
        \cmidrule(lr){2-7}
        & MBJP & CodeT5-220M         & 10.45 & 20.90 & 8.19 & 35.52 \\
        &      & \mainname-220M      & \textbf{11.94} & \textbf{28.36} & \textbf{7.94} & \textbf{1.53} \\
\midrule
\multirow{8}{*}{2B}
        & SuFu & T5Gemma2-2B         & 25.86 & 37.93 & 6.66 & 71.21 \\
        &      & \mainname-2B        & \textbf{36.21} & \textbf{48.28} & \textbf{5.53} & \textbf{0.00} \\
        \cmidrule(lr){2-7}
        & MBJP & T5Gemma2-2B         & 13.43 & 32.84 & 7.46 & 29.55 \\
        &      & \mainname-2B        & \textbf{25.37} & \textbf{43.28} & \textbf{6.36} & \textbf{0.45} \\
        \cmidrule(lr){2-7}
        & HumanEval & T5Gemma2-2B & 31.25 & 37.50 & 6.50 & 24.38 \\
        &      & \mainname-2B        & \textbf{50.00} & \textbf{56.25} & \textbf{4.75} & \textbf{1.30} \\
        \cmidrule(lr){2-7}
        & GFG & T5Gemma2-2B & 19.42 & 33.01 & 7.10 & 13.69 \\
        &      & \mainname-2B        & \textbf{30.10} & \textbf{46.60} & \textbf{5.81} & \textbf{2.75} \\
\bottomrule
\end{tabular}
\begin{tablenotes}
\footnotesize
\item \updir\, Higher is better; \downdir\, Lower is better; metrics in \autoref{sec:evaluation-metrics}.
\end{tablenotes}
\end{threeparttable}
\end{table}

\mainname improves code generation accuracy at both scales and on every benchmark: pass@1 rises from 24.14\% to 37.93\% (SuFu) and from 10.45\% to 11.94\% (MBJP) at 220M, and from 25.86\% to 36.21\% (SuFu), 13.43\% to 25.37\% (MBJP), 31.25\% to 50.00\% (HumanEval-Java), and 19.42\% to 30.10\% (TransCoder-GFG) at 2B; pass@10 improves by 7.46--20.69 points on the same six benchmarks.
The largest gains are on SuFu and HumanEval-Java, where the type constraints remove most uncompilable candidates while the base models often still produce them.
At 220M, the Java baseline has a lower CER than the SuFu baseline (35.52\% versus 83.10\%), leaving fewer compilation errors for type constraints to address. On MBJP, \mainname{} raises pass@10 from 20.90\% to 28.36\% and reduces CER to 1.53\%; the smaller pass@1 gain indicates that finding a correct candidate does not always make it the top-ranked output. The failure analysis in Appendix~\ref{appendix-failure} identifies ranking errors and compilable but incorrect programs among the remaining failures.
\mainname{} also drops CER to 0.00\% on SuFu and 0.45\% / 1.30\% / 2.75\% on the three Java benchmarks; the small residual Java CER comes from static-analysis errors the type system does not capture.

To contextualize absolute performance, \autoref{tab:decoder-only-compare} reports three larger decoder-only models (MiMo-7B, Qwen3-14B, Qwen3-30B-A3B), prompted with the same task prompt in all conditions, alongside our two fine-tuned 2B encoder-decoder models.
Each prompted model produces one greedy candidate per task: the zero-shot condition uses only the task prompt, and the few-shot condition prepends three solved examples for the Java benchmarks and six examples for SuFu.

On SuFu the comparison favors \mainname{}: none of the three solves a single task zero-shot, and few-shot prompting reaches at most 8 of 58, whereas \mainname-2B solves 21; the low-resource language is where the type-guided representation delivers its advantage.
On Java, given the task prompt without examples, the three larger models solve 34--42 of 67 MBJP tasks, 3--14 of 16 HumanEval-Java tasks, and 38--52 of 103 TransCoder-GFG tasks at pass@1. \mainname-2B solves 17/67, 8/16, and 31/103, respectively. These results provide context for performance on Java under the respective generation settings. Moreover, few-shot prompting on Java brings little or no gain over zero-shot for these models, which suggests that they have seen this language extensively and may even have been trained on these benchmarks directly; we therefore flag a possible data-leakage risk for the Java comparison.
In any case, relative to its fine-tuned baseline T5Gemma2-2B, \mainname-2B delivers a consistent improvement on all four benchmarks.

\begin{table}[ht]
\centering
\begin{threeparttable}
\caption{Pass@1 Comparison with Larger Decoder-Only Models}
\label{tab:decoder-only-compare}
\begin{tabular}{l cc cc cc cc}
\toprule
& \multicolumn{2}{c}{MBJP (67)} & \multicolumn{2}{c}{HumanEval-Java (16)} & \multicolumn{2}{c}{TransCoder-GFG (103)} & \multicolumn{2}{c}{SuFu (58)} \\
\cmidrule(lr){2-3} \cmidrule(lr){4-5} \cmidrule(lr){6-7} \cmidrule(lr){8-9}
Model & 0-shot & few-shot & 0-shot & few-shot & 0-shot & few-shot & 0-shot & few-shot \\
\midrule
MiMo-7B             & 34 & 33 & 3  & 10 & 38 & 43 & 0 & 6 \\
Qwen3-14B           & 42 & 41 & 14 & 11 & 50 & 49 & 0 & 8 \\
Qwen3-30B-A3B       & 42 & 45 & 13 & 14 & 52 & 52 & 0 & 6 \\
\midrule
T5Gemma2-2B   & \multicolumn{2}{c}{9}  & \multicolumn{2}{c}{5}  & \multicolumn{2}{c}{20} & \multicolumn{2}{c}{15} \\
\mainname-2B   & \multicolumn{2}{c}{17} & \multicolumn{2}{c}{8}  & \multicolumn{2}{c}{31} & \multicolumn{2}{c}{21} \\
\bottomrule
\end{tabular}
\begin{tablenotes}
\footnotesize
\item Cell entries are solved tasks at pass@1 for all models; the 0-shot condition uses the same prompt as the two fine-tuned models below.
\end{tablenotes}
\end{threeparttable}
\end{table}

\subsubsection{RQ2: Contributions of Individual Components in \mainname}
We conduct an ablation study on the SuFu benchmark and the \mainname-220M model to quantify the contributions of each component in \mainname. Starting from the base case where \mainname synthesizes decision sequences with no checks, system components (syntactic pruning, type pruning, and dynamic typing context, as introduced in \autoref{sec:implementation}) are incrementally added, and the impact on performance metrics is evaluated at each stage. 
\autoref{tab:ablation-combined} summarizes the results.

\begin{table}[ht]
\centering
\setlength{\tabcolsep}{4pt}
\begin{threeparttable}
\caption{Sequential Ablation Study: Performance with Incremental Component Additions}
\label{tab:ablation-combined}
\begin{tabular}{
    l
    S[table-format=2.2]@{\hspace{1em}}S[table-format=+2.2]
    S[table-format=2.2]@{\hspace{1em}}S[table-format=+2.2]
    S[table-format=1.2]@{\hspace{1em}}S[table-format=+1.2]
    S[table-format=2.2]@{\hspace{1em}}S[table-format=+2.2]
    @{\hspace{0.5em}}S[table-format=2.2]@{\hspace{0.5em}}S[table-format=2.1]
}
\toprule
\multicolumn{1}{l}{Components} &
\multicolumn{2}{c}{pass@1 (\%)\tnote{\updir}} &
\multicolumn{2}{c}{pass@10 (\%)\tnote{\updir}} &
\multicolumn{2}{c}{FSP\tnote{\downdir}} &
\multicolumn{2}{c}{CER (\%)\tnote{\downdir}} &
\multicolumn{1}{c}{Time (s)\tnote{\S}} &
\multicolumn{1}{c}{ms/step\tnote{\S}} \\
\cmidrule(lr){2-3} \cmidrule(lr){4-5} \cmidrule(lr){6-7} \cmidrule{8-9} \cmidrule{10-11}
 & {Score} & {$\Delta$} & {Score} & {$\Delta$} & {Score} & {$\Delta$} & {Score} & {$\Delta$} & {} & {} \\
\midrule
Base (no checks)         & 24.14 & 0.00
                        & 36.21 & 0.00
                        & 6.78  & 0.00
                        & 74.26 & 0.00
                        & 5.70 & 31.1 \\
+ Syntactic Pruning         & 27.59 & +3.45
                        & 36.21 & 0.00
                        & 6.71  & -0.07
                        & 72.13 & -2.13
                        & 5.34 & 32.4 \\
+ Type Pruning            & \textbf{37.93} & \textbf{+13.79}
                        & 43.10 & +6.89
                        & 5.79  & -0.99
                        & \textbf{0.00} & \textbf{-74.26}
                        & 11.07 & 43.7 \\
+ Dynamic Typing Context       & \textbf{37.93} & \textbf{+13.79}
                        & \textbf{53.45} & \textbf{+17.24}
                        & \textbf{5.03}  & \textbf{-1.75}
                        & \textbf{0.00}  & \textbf{-74.26}
                        & 15.62 & 59.6 \\
\bottomrule
\end{tabular}
\begin{tablenotes}
\footnotesize
\item[\,] \textit{Score}: absolute value for each setting; $\Delta$: improvement relative to the base setting.
\item[\S] \textit{Time}: average per-task wall time; \textit{ms/step}: average wall time per decoding step.
\end{tablenotes}
\end{threeparttable}
\end{table}

\paragraph{Functional contribution of each component.}
The four metric columns of \autoref{tab:ablation-combined} show which component contributes which improvement.
Syntactic pruning alone gives a moderate pass@1 gain ($+3.45$); the turning point is type pruning, which lifts pass@1 by $+13.79$ and eliminates compilation errors outright (CER $74.26 \to 0.00$).
Dynamic typing context leaves pass@1 and CER unchanged but improves pass@10 ($+17.24$ over the base) and FSP: it ranks correct candidates earlier rather than solving new tasks, which demonstrates the \textbf{Context Locality} property of~\autoref{section:intro}.

\paragraph{How much the constraints prune.}
To measure the pruning strength, we instrument the constrained decoder on the SuFu tasks (beam 10, 330-step budget\footnote{The longest decision sequence in the benchmark has 330 steps, so the budget suffices to derive every reference program.}): syntactic pruning removes 59.1\% of all vocabulary proposals, and type pruning a further 8.2\%.
The lower rate is expected: type checking is more expensive, so it is applied only to the high-probability candidates about to enter the beam.
Such candidates are rarely ill-typed, yet pruning them matters, for they would otherwise expand within the beam and crowd out well-typed derivations.
Type checking during search is provisional (only completed AST fields are checked); the whole program is checked once when the derivation finishes, and this final check rejects 21.2\% of the completed candidates.

\paragraph{Whether the constraints over-prune.}
Completeness guarantees that a well-typed program has a derivation, but it does not guarantee that the model ranks a correct prefix within a fixed beam or provide a bound on the required beam size.
A finite beam can discard all correct prefixes, and the retained hypotheses may later have no valid continuation, leaving the active beam empty.
Training on correct decision sequences is intended to rank correct prefixes highly, but cannot rule out this case.
We measure the effective branching factor as the number of active hypotheses present at each decoding step after pruning.
With beam 10, this number averages 7.97.
It is lower than 10 because completed programs are moved to the output set, while end-of-sequence and rejected candidates do not fill active slots; the top-20 candidate list can therefore be exhausted before all slots are filled.
In our experiments, however, the active beam never became empty.


\paragraph{Runtime cost.}
Each component increases the per-step cost ($31.1 \to 32.4 \to 43.7 \to 59.6$\,ms).
Syntactic pruning adds little (roughly $+4\%$); type pruning and the dynamic typing context cost considerably more, as the former invokes the type checker at every step and the latter additionally re-encodes the evolving synthesis goal.
The per-task totals move by $-6.38\%$ / $+94.00\%$ / $+173.88\%$ ($2.74\times$ base).
Syntactic pruning even lowers the per-task total: with little per-step overhead, it prunes ill-formed decisions early, so beams stay on well-formed derivations that finish sooner (164.6 versus 183.2 mean steps per task), and the $10\%$ fewer steps outweigh the per-step increase.

\subsubsection{RQ3: Comparison with Reliability-Oriented Generation Methods}
We compare \mainname{} with methods that pursue correct code by constraining or repairing the output of a trained model.
On MBJP, four such methods are applied on top of the same frozen T5Gemma2-2B checkpoint selected in RQ1, adding their own pruning or repair mechanisms on top: \emph{rejection sampling}, \emph{SynCode}~\cite{ugare2024syncodellmgenerationgrammar}, \emph{Repilot}(\emph{Copiloting the Copilots})~\cite{DBLP:conf/sigsoft/0003X023}, and \emph{iterative compiler repair}.
All methods are evaluated on the same 67 tasks with 10 candidates each.
Because rejection sampling and repair need diverse draws rather than beam output, the ordinary control uses the temperature-sampling decoding protocol (temperature 0.8, top-$p$ 0.95), so its scores (14.93/23.88) differ from the beam-10 row of \autoref{tab:model-results} (13.43/32.84); the \mainname-2B row is identical in both tables.

Rejection sampling repeatedly draws candidates and keeps those that pass Java compilation, for up to ten ten-draw rounds (100 candidates) per task.
SynCode constrains decoding itself: at each step it masks the tokens that would violate a given grammar, so every output is syntactically well-formed. It is not originally implemented for Java, so we adapt it to the grammar of our Java subset.
Repilot pairs the LM with the Eclipse JDT completion engine: the partial program is maintained as a JDT document, and at each decoding position the engine is queried for the continuations it can legally complete, so that LM-proposed tokens inadmissible at that position are pruned.
Iterative compiler repair feeds \texttt{javac} diagnostics back to the model and regenerates the candidate for at most two rounds; the repair agent is the un-fine-tuned T5Gemma2-2B, since the fine-tuned model never learned to consume diagnostics.

\begin{table}[ht]
\centering
\begin{threeparttable}
\caption{Comparison of \mainname with Reliability-Oriented Generation Methods on Java}
\label{tab:model-compare-sufu-java}
\begin{tabular}{
    l
    S[table-format=2.2]
    S[table-format=2.2]
    S[table-format=2.2]
}
\toprule
Method & {pass@1 (\%)\tnote{\updir}} & {pass@10 (\%)\tnote{\updir}} & {CER (\%)\tnote{\downdir}} \\
\midrule
T5Gemma2-2B           & 14.93 & 23.88 & 26.12 \\
\quad + Rejection Sampling & 16.42 & 28.36 & 0.00 \\
\quad + SynCode (adapt., compile-safe) & 14.93 & 23.88 & 21.34 \\
\quad + Repilot/JDT       & 14.93 & 23.88 & 25.97 \\
\quad + Iterative compiler repair & 14.93 & 23.88 & 17.31 \\
\mainname-2B          & \textbf{25.37} & \textbf{43.28} & \textbf{0.45} \\
\bottomrule
\end{tabular}
\begin{tablenotes}
\footnotesize
\item \updir\,/\downdir\, Higher/lower is better.
\end{tablenotes}
\end{threeparttable}
\end{table}

None of the four external baselines significantly improves functional correctness over the ordinary control:
rejection sampling adds three solved tasks at pass@10 (19 vs.\ 16 of 67) and one at pass@1 (11 vs.\ 10; McNemar $p=1.00$ and $p=0.25$);
SynCode, Repilot, and iterative repair leave the solved sets exactly unchanged (10/67 and 16/67);
Repilot's pruning rarely fires: following the original Repilot implementation, punctuation and keyword tokens are accepted without querying the engine, and the completion list at an incomplete position is permissive, so compilation errors barely move (26.12\% $\to$ 25.97\%);
iterative repair repairs 59 candidates' compilability (compile errors 175$\to$116 of 670) yet leaves the solved sets unchanged.
\mainname-2B solves 17/67 at pass@1 and 29/67 at pass@10 with 3 of 670 candidates failing to compile, the best of all compared methods.
External reliability mechanisms (post-hoc filtering, token-level grammar masking, completion-engine repair, compiler-feedback regeneration) never change the model's own preference and distribution over output tokens, whereas \mainname{} explicitly encodes the typing rules into the model and helps it better learn the program distribution, demonstrating the importance of the \textbf{Type Explicitness} property in~\autoref{section:intro}.

\subsubsection{RQ4: Comparison with Separated Type-Code Generation}
We compare \mainname with approaches that separate type reasoning and code generation.
Both variants fine-tune CodeT5-220M to emit a single sequence that concatenates the two parts, differing only in their order: \emph{Type-First} generates the sequence of typing rules (forming the type derivation tree) first and then the code on its basis, whereas \emph{Code-First} generates the code first and then the typing rules as its type reasoning.
All settings are evaluated on the same SuFu and MBJP test sets with 10 candidates each.

\begin{table}[ht]
\centering
\begin{threeparttable}
\caption{Comparison: Separated vs. Integrated Type-Code Generation}
\label{tab:sequential-compare}
\begin{tabular}{
    l
    l
    S[table-format=2.2]
    S[table-format=2.2]
    S[table-format=2.2]
    S[table-format=1.2]
    S[table-format=3.0]
}
\toprule
Language & Model & {pass@1 (\%)\tnote{\updir}} & {pass@10 (\%)\tnote{\updir}} & {CER (\%)\tnote{\downdir}} & {FSP\tnote{\downdir}} & {Avg Tokens\tnote{\downdir}} \\
\midrule
\multirow{4}{*}{SuFu}
        & CodeT5-220M                & 24.14 & 32.76 & 83.10 & 7.03 & \textbf{405} \\
        & CodeT5-220M + Type-First   & 31.03 & 36.21 & 64.31 & 6.57 & 752 \\
        & CodeT5-220M + Code-First   & 29.31 & 39.66 & 56.72 & 6.28 & 752 \\
        & \mainname-220M             & \textbf{37.93} & \textbf{53.45} & \textbf{0.00} & \textbf{5.03} & 410 \\
\midrule
\multirow{4}{*}{Java}
        & CodeT5-220M                & 10.45 & 20.90 & 35.52 & 8.19 & 142 \\
        & CodeT5-220M + Type-First   & \textbf{11.94} & 26.87 & 37.61 & \textbf{7.36} & 247 \\
        & CodeT5-220M + Code-First   & 8.96 & 22.39 & 31.49 & 8.21 & 247 \\
        & \mainname-220M             & \textbf{11.94} & \textbf{28.36} & \textbf{1.53} & 7.94 & \textbf{129} \\
\bottomrule
\end{tabular}
\begin{tablenotes}
\footnotesize
\item \updir\,/\downdir\, Higher/lower is better.
\item \emph{Avg Tokens}: length of the generated sequence in each model's own vocabulary, excluding comments.
\end{tablenotes}
\end{threeparttable}
\end{table}

Adding separated type reasoning improves over the ordinary trained baseline on both languages (e.g., Type-First 31.03 vs.\ 24.14 pass@1 on SuFu), and unifying the two into one process improves further: both separated variants trail \mainname{} on pass@10 and CER (on Java pass@1, Type-First merely ties \mainname{}).
Because neither variant constrains decoding, their CER remains high (56.72--64.31 on SuFu), whereas \mainname{} outputs only type-valid programs.
In terms of token length, fully unifying type reasoning and code is far more economical than outputting them separately: the unified decision sequence (410 on SuFu, 129 on Java) saves 342 and 118 tokens respectively against appending the typing rules to plain code tokens (752 and 247).
It is even comparable to, or slightly shorter than, the plain code tokens alone (405 on SuFu, 142 on Java): a type derivation rule absorbs the redundant syntactic elements of the code (such as equals signs and semicolons) into a single rule token, so our representation adds little and sometimes removes tokens.
These results support the \textbf{Derivation Vicinality} paradigm of \mainname, that interleaving each program fragment with its type derivation is more accurate and more efficient than generating them separately.

\subsection{Limitations and Future Work} \label{sec:limitations}
\paragraph{Dependence on the encoder-decoder architecture.}
The current architecture is designed for encoder-decoder models: at each synthesis step, an encoder processes the evolving synthesis goal to provide the dynamic typing context.
The representation itself carries over to a decoder-only model unchanged; what changes is the architecture, and with it the cost of the dynamic typing context.
In a decoder-only model, supplying the context would mean prefilling the current typing context after the derivation prefix at every step and removing it again before the next step, so its key-value entries are recomputed at every step and never reused by later steps, and the computation is awkward to schedule.
We leave this adaptation to future work.

\paragraph{Dynamically typed languages.}
The current design of \mainname relies on a type system in which the type derivation of a program can be computed automatically from the program text, so that existing code can be converted into our decision-sequence format.
Dynamically typed languages such as Python still have typing rules to a degree, but the type derivation cannot be read off the program text alone, since types depend on control flow and runtime information.
As a result, decision sequences cannot be automatically extracted from existing programs, which breaks the \textbf{Data Usability} property of~\autoref{section:intro} and with it the training-data pipeline.
Parsing and adapting such languages is left to future work.

\paragraph{Type-system expressiveness.}
Both evaluated type systems are relatively simple compared with the mainstream languages that have evolved over decades, such as C++: the implemented Java subset omits subtyping hierarchies, overload resolution, and mutable-state-related features, and SuFu's type system targets fusion optimization.
The framework itself is not confined to such systems: typing rules are formulated as constrained Horn clauses (Sec.~\ref{sec:meta}), which can in principle describe more complex type systems and their languages.
For features that can be represented with first-order terms, such as many forms of subtyping, overloading, effects, and type classes, we need to implement the typing rules and solvers and then measure their cost.
The current solver cannot handle a typing rule that requires higher-order unification, as discussed in Sec.~\ref{sec:meta}.
Our experiments do not test either type of extension. We leave them for future work.
